\chapter{Transform Coding}

\section{Table of Contents}

\begin{itemize}
  \item Core Idea
  \item Key Concepts
  \item Theory and Formulas
  \item Visual Schemes
  \item Examples
\end{itemize}

\section{Core Idea}

\textbf{Transform coding} improves compression by converting correlated samples into coefficients with very unequal importance. Important coefficients receive more bits; negligible coefficients receive few bits or become zero.

Scalar quantization of natural signals is weak because samples are correlated and often have similar variance. Transform coding applies a reversible orthogonal transform $Y=TX$ so that:

\begin{itemize}
  \item most energy is concentrated in few coefficients,
  \item coefficient variances become diverse,
  \item high-variance coefficients get more bits,
  \item low-variance coefficients can be heavily quantized.
\end{itemize}

\begin{calloutbox}{red!70!black}{Transform coding principle}
Compression gain does not come from losing information in the transform. It comes from making coefficient variances uneven, then allocating quantization bits where they matter.
\end{calloutbox}

\section{Key Concepts}

\subsection{Block Coding and Rate Allocation}

Consider a block:

\[
X=[X_1,X_2,\ldots,X_M]^T
\]

For component $k$:

\begin{itemize}
  \item $\sigma_k^2=\mathrm{Var}(X_k)$,
  \item $c_k$ is the shape factor,
  \item $R_k$ is the rate assigned to component $k$.
\end{itemize}

With optimal scalar quantization per component:

\[
D_k=c_k\sigma_k^2 2^{-2R_k}
\]

Global MSE per component:

\[
\mathcal{D}
=\frac{1}{M}\mathbb{E}\|X-Q(X)\|^2
=\frac{1}{M}\sum_{k=0}^{M-1}c_k\sigma_k^2 2^{-2R_k}
\]

Rate allocation problem:

\[
\min_R \mathcal{D}(R)
\quad\text{s.t.}\quad
\sum_{k=0}^{M-1}R_k\leq R_{\text{Tot}}
\]

\subsection{Huang-Schulteiss Formula}

\begin{calloutbox}{red!70!black}{Optimal bit allocation}
\[
R_k^*
=
\frac{R_{\text{Tot}}}{M}
+
\frac{1}{2}\log_2
\left(
\frac{c_k\sigma_k^2}
{c_{\text{GM}}\sigma_{\text{GM}}^2}
\right)
\]

where:

\[
c_{\text{GM}}=\left(\prod_k c_k\right)^{1/M},
\qquad
\sigma_{\text{GM}}^2=\left(\prod_k\sigma_k^2\right)^{1/M}
\]

More variance means more bits. At optimum, active components have equal distortion.
\end{calloutbox}

Optimal distortion:

\[
\mathcal{D}^*
=c_{\text{GM}}\sigma_{\text{GM}}^2 2^{-2\bar{R}},
\qquad
\bar{R}=\frac{R_{\text{Tot}}}{M}
\]

For Gaussian variables, $c_k=c_N$:

\[
R_k^*=\bar{R}+\frac{1}{2}\log_2\frac{\sigma_k^2}{\sigma_{\text{GM}}^2}
\]

If all components are identically distributed:

\[
R_k^*=\bar{R},
\qquad
\mathcal{D}=c_X\sigma_X^2 2^{-2\bar{R}}
\]

so block coding alone gives no gain. Gain requires \textbf{variance diversity}.

\subsection{Orthogonal Transforms}

An orthogonal transform satisfies:

\[
T^{-1}=T^T
\]

Forward and inverse:

\[
Y=TX,\qquad X=T^TY
\]

Orthogonal transforms preserve Euclidean norm:

\[
\|TX\|^2=X^T(T^TT)X=\|X\|^2
\]

\begin{calloutbox}{red!70!black}{Distortion preservation}
\[
\mathcal{D}_Y
=
\frac{1}{M}\mathbb{E}\|Y-\hat{Y}\|^2
=
\frac{1}{M}\mathbb{E}\|T(X-\hat{X})\|^2
=
\mathcal{D}_X
\]

Therefore, quantization distortion can be analyzed in transform domain without changing MSE.
\end{calloutbox}

\subsection{Coding Gain}

For i.d. Gaussian samples, direct PCM gives:

\[
D_{\text{PCM}}=c_N\sigma_X^2 2^{-2\bar{R}}
\]

After orthogonal transform and optimal allocation:

\[
\mathcal{D}_T=c_N\sigma_{\text{GM},Y}^2 2^{-2\bar{R}}
\]

Orthogonal transforms preserve average variance:

\[
\sigma_{\text{AM},Y}^2
=
\frac{1}{M}\mathbb{E}\|Y\|^2
=
\frac{1}{M}\mathbb{E}\|X\|^2
=\sigma_X^2
\]

\begin{calloutbox}{red!70!black}{Coding gain}
\[
G_T
=
\frac{D_{\text{PCM}}}{\mathcal{D}_T}
=
\frac{\sigma_{\text{AM},Y}^2}
{\sigma_{\text{GM},Y}^2}
\geq 1
\]

Since arithmetic mean is at least geometric mean, gain increases when transform coefficients have very unequal variances.
\end{calloutbox}

\subsection{Practical Rate Allocation}

The HS formula can produce negative or fractional rates. Practical schemes:

\begin{center}\fitwidth{%
\begin{tabular}{|l|l|l|}
\hline
\textbf{Method} & \textbf{Idea} & \textbf{Use} \\ \hline
\textbf{Modified HS} & remove negative-rate components, recompute, floor, distribute residual bits & continuous optimum adapted to integer rates \\ \hline
\textbf{Greedy} & start at zero bits, repeatedly add one bit to component with largest current distortion & same result, simpler for small total rate \\ \hline
\end{tabular}}
\end{center}

Greedy update after adding one bit:

\[
D_k \leftarrow \frac{D_k}{4}
\]

because:

\[
2^{-2(R_k+1)}=\frac{1}{4}2^{-2R_k}
\]

\subsection{KLT}

Let $X$ be zero-mean with covariance/correlation matrix:

\[
R_X=\mathbb{E}[XX^T]
\]

\begin{calloutbox}{red!70!black}{Karhunen-Loeve Transform}
$T_{\text{KLT}}$ is the orthonormal matrix whose rows are eigenvectors of $R_X$:

\[
T_{\text{KLT}}=[u_1\ u_2\ \cdots\ u_M]^T,
\qquad
Y=T_{\text{KLT}}X
\]

It decorrelates coefficients:

\[
\mathbb{E}[Y_iY_j]=\lambda_i\delta_{ij}
\]

It gives best energy compaction among orthogonal transforms and maximizes coding gain for Gaussian sources.
\end{calloutbox}

KLT limitations:

\begin{itemize}
  \item data-dependent basis,
  \item expensive eigenvector computation,
  \item transform matrix must be known at decoder,
  \item stationarity assumptions may fail for images.
\end{itemize}

For strongly correlated AR(1)-like sources, fixed frequency transforms such as DCT approach KLT performance with far lower cost.

\subsection{DFT}

1D orthonormal DFT:

\[
y[k]
=
\frac{1}{\sqrt{M}}\sum_{n=1}^{M}x[n]e^{-j\frac{2\pi}{M}kn},
\qquad
k=1,\ldots,M
\]

DFT matrix:

\[
(T_{\text{DFT}})_{k,n}
=
\frac{1}{\sqrt{M}}W_M^{kn},
\qquad
W_M=e^{-j2\pi/M}
\]

2D separable transform:

\[
Y=TXT^T
\]

2D basis functions:

\[
B_{k,\ell}(n,m)
=
\frac{1}{N}e^{j\frac{2\pi}{N}(kn+\ell m)}
\]

DFT is not used directly in image compression because it implicitly periodizes blocks. If block boundaries do not match, artificial jumps appear, causing \textbf{spectral leakage} and too many high-frequency coefficients.

\subsection{DCT}

DCT solves DFT leakage by using symmetric extension before periodization. This removes boundary discontinuities and produces real coefficients.

\begin{calloutbox}{red!70!black}{DCT matrix}
\[
(T_{\text{DCT}})_{k,n}
=
\begin{cases}
\frac{1}{\sqrt{M}} & k=0\\
\sqrt{\frac{2}{M}}
\cos\frac{(2n+1)k\pi}{2M} & k>0
\end{cases}
\]

for $k,n=0,\ldots,M-1$.

DCT gives real coefficients, avoids negative frequencies, reduces leakage, and is the standard transform for compression.
\end{calloutbox}

2D DCT is separable:

\[
Y=TXT^T
\]

In image blocks:

\begin{itemize}
  \item top-left coefficient is \textbf{DC},
  \item other coefficients are \textbf{AC} spatial frequencies,
  \item energy usually decays toward bottom-right,
  \item smooth blocks become very sparse,
  \item textured or edge blocks keep more coefficients.
\end{itemize}

\subsection{JPEG Overview}

Baseline JPEG is a block-based DCT image codec. Standard decoder is specified; encoders can choose implementation details.

Pipeline:

\begin{enumerate}
  \item RGB to YCbCr color conversion.
  \item Optional chroma subsampling.
  \item Split each component into $8x8$ blocks.
  \item Subtract 128 from samples.
  \item Apply 2D DCT.
  \item Quantize DCT coefficients.
  \item Entropy-code DC and AC coefficients.
\end{enumerate}

\subsection{JPEG Quantization}

Per coefficient:

\[
\tilde{c}_{i,j}
=
\mathrm{round}\left(\frac{c_{i,j}}{q_{i,j}}\right)
\]

De-quantization:

\[
\hat{c}_{i,j}=\tilde{c}_{i,j}q_{i,j}
\]

Luminance standard table:

\[
q^*=
\begin{pmatrix}
16&11&10&16&24&40&51&61\\
12&12&14&19&26&58&60&55\\
14&13&16&24&40&57&69&56\\
14&17&22&29&51&87&81&61\\
18&22&37&56&68&109&103&77\\
24&35&55&64&81&104&111&90\\
49&63&78&87&101&121&120&100\\
72&92&95&98&112&100&103&99
\end{pmatrix}
\]

Values grow toward high frequencies, so high-frequency coefficients are quantized more coarsely.

Quality factor scaling:

\[
S_F=
\begin{cases}
5000/Q & 1\leq Q\leq50\\
200-2Q & 50<Q\leq99\\
1 & Q=100
\end{cases}
\qquad
q=\frac{S_F}{100}q^*
\]

\subsection{JPEG Entropy Coding}

After quantization, most high-frequency coefficients are zero.

\begin{enumerate}
  \item \textbf{Zig-zag scan} orders coefficients from low to high spatial frequency.
  \item \textbf{DC coefficient} is coded by DPCM:
\end{enumerate}

\[
DC_P=DC_n-DC_{n-1}
\]

\begin{enumerate}
  \item \textbf{AC coefficients} are coded by run-length plus category/amplitude.
\end{enumerate}

DC category:

\[
k=\lceil\log_2(|DC_P|+1)\rceil
\]

AC symbols:

\begin{itemize}
  \item run-length = number of zeros before next non-zero coefficient,
  \item category = amplitude magnitude class,
  \item amplitude = signed value bits.
\end{itemize}

Special AC symbols:

\begin{center}\fitwidth{%
\begin{tabular}{|l|l|}
\hline
\textbf{Symbol} & \textbf{Meaning} \\ \hline
$(15,0)$ & zero-run, 16 consecutive zeros \\ \hline
$(0,0)$ & end of block \\ \hline
\end{tabular}}
\end{center}

\subsection{JPEG File Structure}

JPEG bitstream hierarchy:

\begin{verbatim}
Frame
|-- Frame header
`-- Scan
    |-- Scan header
    `-- Segment
        |-- Segment header
        `-- 8x8 blocks
\end{verbatim}

JFIF stores basic interchange metadata. Exif stores camera-oriented metadata such as exposure, ISO, GPS, date, thumbnail.

\section{Theory and Formulas}

\subsection{HS Derivation Sketch}

Lagrangian:

\[
J(R,\lambda)
=
\frac{1}{M}\sum_{k=0}^{M-1}c_k\sigma_k^2 2^{-2R_k}
+
\lambda\left(\sum_{k=0}^{M-1}R_k-R_{\text{Tot}}\right)
\]

Stationarity:

\[
\frac{\partial J}{\partial R_k}=0
\]

gives:

\[
2^{-2R_k^*}
=
\frac{M\lambda}{2\ln2}
\frac{1}{c_k\sigma_k^2}
\]

Therefore:

\[
R_k^*
=
\lambda'
+
\frac{1}{2}\log_2(c_k\sigma_k^2)
\]

The constraint $\sum_kR_k^*=R_{\text{Tot}}$ determines:

\[
\lambda'
=
\bar{R}
-
\frac{1}{2}\log_2(c_{\text{GM}}\sigma_{\text{GM}}^2)
\]

\subsection{AM-GM and Coding Gain}

For non-negative variances:

\[
\sigma_{\text{GM}}^2
=
\left(\prod_k\sigma_k^2\right)^{1/M}
\leq
\frac{1}{M}\sum_k\sigma_k^2
=
\sigma_{\text{AM}}^2
\]

Equality holds only if all variances are equal. Therefore a transform improves coding only when it makes coefficient variances unequal.

\subsection{Toy Rotation Example}

For two highly correlated variables, independent scalar quantization wastes bits because the joint PDF lies on a thin diagonal region.

Use a 45-degree orthogonal rotation:

\[
\begin{bmatrix}
Y_1\\Y_2
\end{bmatrix}
=
\frac{1}{\sqrt{2}}
\begin{bmatrix}
1&-1\\
1&1
\end{bmatrix}
\begin{bmatrix}
X_1\\X_2
\end{bmatrix}
\]

After transform:

\[
\sigma_{Y_1}^2=\frac{\Delta_1^2}{12}=2\sigma^2,
\qquad
\sigma_{Y_2}^2=\frac{\Delta_2^2}{12}\ll\sigma^2
\]

So $Y_2$ can be ignored or heavily quantized, while most bits go to $Y_1$.

\subsection{KLT Energy Compaction}

For any $N<M$ and any other orthogonal transform $Z=TX$:

\[
\sum_{i=1}^{N}\mathbb{E}[Y_i^2]
\geq
\sum_{i=1}^{N}\mathbb{E}[Z_i^2]
\]

In multispectral imaging example:

\begin{center}\fitwidth{%
\begin{tabular}{|l|l|}
\hline
\textbf{KLT band} & \textbf{Energy} \\ \hline
1 & 86.37\% \\ \hline
2 & 7.65\% \\ \hline
3 & 4.16\% \\ \hline
4 & 1.67\% \\ \hline
5 & 0.12\% \\ \hline
6 & 0.03\% \\ \hline
\end{tabular}}
\end{center}

First KLT band carries most energy; later bands can be quantized aggressively.

\subsection{JPEG Rate-Distortion Notes}

JPEG's quality factor controls quantization but not exact bitrate. Bitrate depends on:

\begin{itemize}
  \item image content,
  \item DCT coefficient sparsity,
  \item quantization table,
  \item DC differences,
  \item AC run lengths,
  \item Huffman tables.
\end{itemize}

This is why JPEG has weak direct rate control compared with modern codecs.

\section{Visual Schemes}

\subsection{Transform Sparsification}

\begin{center}\includegraphics[width=0.917\linewidth]{ch04/dft-sparsification.png}\end{center}
A pure tone is dense in sample domain but sparse in Fourier domain, showing why transforms can help compression.

\begin{center}\includegraphics[width=0.7\linewidth]{ch04/correlated-variables-before-transform.png}\end{center}
Correlated variables have an elongated joint density; independent scalar quantization does not exploit this geometry.

\begin{center}\includegraphics[width=0.7\linewidth]{ch04/correlated-variables-after-transform.png}\end{center}
A rotation aligns the axes with high and low variance directions, making bit allocation effective.

\subsection{KLT, DFT, and DCT}

\begin{center}\includegraphics[width=\linewidth]{ch04/klt-decorrelation.png}\end{center}
KLT aligns the coordinate system with covariance eigenvectors, decorrelating coefficients and compacting energy.

\begin{center}\includegraphics[width=\linewidth]{ch04/dft-natural-image-spectrum.png}\end{center}
Natural-image DFT magnitude concentrates near low frequencies but shows leakage from implicit periodic boundaries.

\begin{center}\includegraphics[width=\linewidth]{ch04/dct-symmetric-extension.png}\end{center}
DCT mirrors the block before periodization, avoiding boundary jumps and reducing spectral leakage.

\begin{center}\includegraphics[width=0.833\linewidth]{ch04/dct-coefficient-energy.png}\end{center}
2D-DCT concentrates most natural-image block energy in top-left low-frequency coefficients.

\begin{center}\includegraphics[width=0.833\linewidth]{ch04/dct-basis-functions.png}\end{center}
Each $8x8$ DCT coefficient measures similarity to one cosine basis pattern.

\subsection{JPEG Pipeline}

\begin{center}\includegraphics[width=\linewidth]{ch04/jpeg-pipeline.png}\end{center}
JPEG combines color conversion, chroma subsampling, block DCT, quantization, and entropy coding.

\begin{center}\footnotesize\itshape Diagram (Mermaid source):\end{center}
\begin{verbatim}
flowchart LR
    A[RGB image] --> B[YCbCr]
    B --> C[Chroma subsampling]
    C --> D[8x8 blocks]
    D --> E[Subtract 128]
    E --> F[2D DCT]
    F --> G[Quantization]
    G --> H[Zig-zag scan]
    H --> I[DC DPCM / AC RLE]
    I --> J[Huffman coding]
    J --> K[JPEG bitstream]
\end{verbatim}

\begin{center}\includegraphics[width=\linewidth]{ch04/jpeg-dct-block-coefficients.png}\end{center}
JPEG blocks usually have large low-frequency DCT coefficients and many near-zero high-frequency coefficients.

\begin{center}\includegraphics[width=0.7\linewidth]{ch04/jpeg-zigzag-scan.png}\end{center}
Zig-zag scan converts a sparse 2D coefficient matrix into a 1D sequence with zeros clustered near the end.

\begin{center}\includegraphics[width=\linewidth]{ch04/jpeg-file-structure.png}\end{center}
JPEG files store frame, scan, segment, table, and block information needed by the decoder.

\subsection{JPEG Artifacts}

\begin{center}\includegraphics[width=0.55\linewidth]{ch04/jpeg-quality-high.png}\end{center}
High-quality JPEG preserves visual detail at higher bitrate.

\begin{center}\includegraphics[width=0.55\linewidth]{ch04/jpeg-quality-low.png}\end{center}
Lower bitrate makes blocking and ringing artifacts visible.

\begin{center}\includegraphics[width=0.55\linewidth]{ch04/jpeg-quality-very-low.png}\end{center}
Very low bitrate strongly reveals $8x8$ block structure and loss of high frequencies.

\section{Examples}

\begin{calloutbox}{green!55!black}{Modified HS allocation}
\textbf{Data:} $\sigma_1^2=1000$, $\sigma_2^2=100$, $\sigma_3^2=50$, $\sigma_4^2=1$, $R_{\text{Tot}}=10$.

HS gives a negative rate for component 4, so it is removed and rates are recomputed over active components.

Final integer allocation:

\[
R=[5,3,2,0]
\]

Distortions:

\[
D_1=1000\cdot2^{-10}\approx0.98
\]

\[
D_2=100\cdot2^{-6}\approx1.56
\]

\[
D_3=50\cdot2^{-4}\approx3.13,\qquad D_4=1
\]

Total distortion is about $6.67$.
\end{calloutbox}

\begin{calloutbox}{green!55!black}{Greedy allocation}
Same data, start with $R_k=0$ and $D_k=\sigma_k^2$.

At each iteration, add one bit to the component with largest current $D_k$ and divide that $D_k$ by 4.

After 10 bits, greedy gives the same allocation:

\[
R=[5,3,2,0]
\]
\end{calloutbox}

\begin{calloutbox}{green!55!black}{Transform coding vs. scalar quantization}
Without transform, scalar quantization of two correlated variables gives:

\[
D=2\sigma^2 2^{-2R}
\]

After rotation, nearly all energy is in $Y_1$ and $Y_2$ has negligible variance. Quantizing only $Y_1$ gives much lower distortion at the same average rate.

\begin{center}\fitwidth{%
\begin{tabular}{|l|l|l|}
\hline
\textbf{Total bits per pair} & \textbf{Average rate} & \textbf{Distortion after transform} \\ \hline
0 & 0 & $2\sigma^2$ \\ \hline
1 & 0.5 bpS & $\sigma^2/2$ \\ \hline
2 & 1 bpS & $\sigma^2/8$ \\ \hline
3 & 1.5 bpS & $\sigma^2/32$ \\ \hline
\end{tabular}}
\end{center}
\end{calloutbox}

\begin{calloutbox}{green!55!black}{JPEG centering and DCT}
Before DCT, JPEG subtracts 128 from each 8-bit sample so values are centered near zero.

Example source block values around 91-184 become approximately -37 to +56 after centering. The DCT then produces 64 coefficients, often with most energy concentrated in low-frequency positions.
\end{calloutbox}

\begin{calloutbox}{green!55!black}{JPEG quantization}
Quantization divides each DCT coefficient by its table entry and rounds:

\[
\tilde{c}_{i,j}=\mathrm{round}(c_{i,j}/q_{i,j})
\]

Example quantized block has only a few non-zero coefficients:

\begin{verbatim}
-2  17   3   1   0   0   0   0
 3   4   0  -1   0   0   0   0
 0   0  -1   0   0   0   0   0
 0   0  -1  -1   0   0   0   0
\end{verbatim}

After de-quantization and inverse DCT, source example has \textbf{MSE = 24.05} and \textbf{PSNR = 34.32 dB}.
\end{calloutbox}

\begin{calloutbox}{green!55!black}{JPEG DC coding}
If:

\[
DC_P=-5
\]

then:

\[
k=\lceil\log_2(|-5|+1)\rceil=3
\]

Category code for $k=3$ is \texttt{100}. Magnitude $5$ is \texttt{101}; negative amplitude flips bits to \texttt{010}.

Final DC code:

\begin{verbatim}
100010
\end{verbatim}
\end{calloutbox}

\begin{calloutbox}{green!55!black}{JPEG AC coding}
AC coding uses tuples:

\[
(\text{run-length},\text{category},\text{amplitude})
\]

A long tail of zero coefficients can be ended with EOB \texttt{(0,0)}. Runs of 16 zeros use ZR \texttt{(15,0)}.

In the source example, the final entropy-coded block has 66 bits for 64 pixels:

\[
66/64\approx1.03\ \text{bpp}
\]
\end{calloutbox}

\begin{calloutbox}{green!55!black}{JPEG quality levels}
\begin{center}\fitwidth{%
\begin{tabular}{|l|l|l|l|}
\hline
\textbf{Quality} & \textbf{Rate} & \textbf{PSNR} & \textbf{Visual effect} \\ \hline
High & 1.02 bpp & 33.92 dB & almost identical \\ \hline
Medium-high & 0.75 bpp & 33.45 dB & minor artifacts \\ \hline
Medium & 0.50 bpp & 32.70 dB & slight blocking \\ \hline
Low & 0.31 bpp & 31.31 dB & visible blocking \\ \hline
Very low & 0.21 bpp & 29.50 dB & strong blocking and ringing \\ \hline
\end{tabular}}
\end{center}

\textbf{Takeaway:} JPEG quality factor controls quantization strength, but exact rate is hard to predict from the factor alone.
\end{calloutbox}
